% Generated draft: requires mathematical and source review.
\documentclass[11pt]{article}
\usepackage[margin=25mm]{geometry}
\usepackage{fontspec}
\setmainfont{Latin Modern Roman}
\usepackage{amsmath,amssymb,mathtools,unicode-math}
\setmathfont{Latin Modern Math}
\usepackage{enumitem,needspace,xurl,hyperref}
\setlength{\parindent}{0pt}
\setlength{\parskip}{5pt}
\setlength{\emergencystretch}{3em}
\newcommand{\N}{\mathbb{N}}
\newcommand{\Z}{\mathbb{Z}}
\newcommand{\Q}{\mathbb{Q}}
\newcommand{\R}{\mathbb{R}}
\newcommand{\C}{\mathbb{C}}
\newcommand{\F}{\mathbb{F}}
\newcommand{\A}{\mathbb{A}}
\newcommand{\PP}{\mathbb{P}}
\newcommand{\E}{\mathbb{E}}
\newcommand{\eps}{\varepsilon}
\newcommand{\dd}{\,\mathrm{d}}
\begin{document}
\section*{UnsolvedMath ID: 30001584}
\textbf{Source problem code:} catalogue entry
\subsection{Definitions and mathematical statement}
Let \(G\) be an ambient group. The supplied record does not specify the intended class of groups \(G\), and the report title alone does not establish any additional hypotheses on \(G\).

For a prime \(p\), write \(\Z/p\Z\) for the cyclic group of order \(p\). Under the standard restricted regular wreath-product interpretation, for a group \(K\) define
\[
(\Z/p\Z)\wr K
=
\left(\bigoplus_{k\in K}\Z/p\Z\right)\rtimes K,
\]
where the direct sum consists of finitely supported functions \(f:K\to\Z/p\Z\), and \(K\) acts by left translation:
\[
(h\cdot f)(k)=f(h^{-1}k)\qquad(h,k\in K).
\]
Here containment in \(G\) is interpreted as the existence of an injective group homomorphism into \(G\). The supplied record does not explicitly specify the wreath-product convention; the displayed definition makes this conventional reading explicit and requires confirmation against the source.

The symbol \(C(G)\) is not defined in the supplied record. Its title and literature assessment associate it with \(\ell^2\)-Betti values, but do not determine which values are included, the coefficient ring, or the operators or spaces used to define them. Retaining \(C(G)\) as the unspecified set intended by the source, the first question is whether
\[
C(G)\not\subseteq\Q
\quad\Longleftrightarrow\quad
\exists\,p\text{ prime}\;\exists\text{ an injective homomorphism }
(\Z/p\Z)\wr\Z\longrightarrow G.
\]
The source then asks, more generally, whether
\[
C(G)\not\subseteq\Q
\quad\Longleftrightarrow\quad
\exists\,p\text{ prime}\;\exists\,K\text{ an infinite finitely generated group}\;
\exists\text{ an injective homomorphism }
(\Z/p\Z)\wr K\longrightarrow G.
\]
In each question, \(C(G)\not\subseteq\Q\) means that some element of \(C(G)\) is not rational. These are two proposed equivalences, not asserted results. A fully specified mathematical statement requires the missing definition of \(C(G)\) and the intended hypotheses on \(G\).
\subsection{Short English statement}
The source asks whether having a non-rational value in the unspecified invariant \(C(G)\) is equivalent to containing a lamplighter subgroup \((\Z/p\Z)\wr\Z\), and asks a second version allowing any infinite finitely generated group \(K\) in place of \(\Z\). The definition of \(C(G)\), the ambient group hypotheses, and the wreath-product convention require source review.
\subsection{Sources}
\begin{itemize}
\item Supplied local UnsolvedMath registry record: permanent ID \(30001584\), exact source problem code \texttt{catalogue entry}, title ``Lamplighter Subgroups and Irrational \(\ell^2\)-Betti Numbers,'' in the Oberwolfach Reports Open Problems collection. The record gives proposed year 2010 and supplies no individual proposer.
\item Supplied source citation: \emph{Actions and Invariants of Residually Finite Groups: Asymptotic Methods} (2011), Oberwolfach Reports, pp.\ 2335--2391, DOI \texttt{10.4171/owr/2010/40}. Original OWR report: \url{https://doi.org/10.4171/owr/2010/40}. The background describes the collection as drawn from workshop reports and their open-problem discussions.
\item Supplied literature reference: Grabowski, \emph{Atiyah problem for lamplighter groups}, \url{https://arxiv.org/abs/1009.0229}.
\item Literature assessment dated exactly \texttt{2026-08-21}: the supplied curation labels the problem open and states that the proposed characterization is not established; it attributes sufficiency in the lamplighter direction to Grabowski's transcendental examples, while stating that these do not prove the converse. This is supplied dated curation, not an independently verified assessment. The record explicitly calls for independent verification of time-sensitive status and cited claims.
\item The registry supplies statement status \texttt{corrected\_verified} and says its mathematical content was checked against the official report, with notation, context, and question form normalized without changing the source-level claim. These are supplied verification assertions; no independent source consultation or mathematical verification is claimed here.
\end{itemize}
\end{document}
